\chapter{2013年新课标I卷（理）}

\section{单选题}

\begin{problem}
已知集合 \(A = \{x \mid x^2 - 2x > 0\}\)，\(B = \{x \mid -\sqrt{5} < x < \sqrt{5}\}\)，则（\quad）

\choices
    {\(A \cap B = \varnothing\)}
    {\(A \cup B = \R\)}
    {\(B \subseteq A\)}
    {\(A \subseteq B\)}
\end{problem}

\begin{answer}
B
\end{answer}

\begin{solution}
由 \(x^2-2x=x(x-2)>0\)，得 \(A=(-\infty,0)\cup(2,+\infty)\)，而 \(B=(-\sqrt{5},\sqrt{5})\). 因为 \(\sqrt{5}>2\)，所以 \([0,2]\subset B\)；当 \(x<0\) 或 \(x>2\) 时，\(x\in A\)，从而 \(A\cup B=\R\). 又 \(-1\in A\cap B\)，故 \(A\cap B\neq\varnothing\)；\(1\in B\) 且 \(1\notin A\)，故 \(B\) 不是 \(A\) 的子集；\(3\in A\) 且 \(3\notin B\)，故 \(A\) 不是 \(B\) 的子集. 因此选 B.
\end{solution}

\begin{problem}
若复数 \(z\) 满足 \(\left(3 - 4\i\right)z = \left|4 + 3\i\right|\)，则 \(z\) 的虚部为（\quad）

\choices
    {\(-4\)}
    {\(-\frac{4}{5}\)}
    {4}
    {\(\frac{4}{5}\)}
\end{problem}

\begin{answer}
D
\end{answer}

\begin{solution}
因为 \(\left|4+3\i\right|=5\)，所以
\(z=\dfrac{5}{3-4\i}=\dfrac{5(3+4\i)}{(3-4\i)(3+4\i)}=\dfrac{3+4\i}{5}\). 因而 \(z\) 的虚部为 \(\dfrac45\)，选 D.
\end{solution}

\begin{problem}
为了解某地区的中小学生视力情况，拟从该地区的中小学生中抽取部分学生进行调查，事先已了解到该地区小学，初中，高中三个学段学生的视力情况有较大差异，而男女生视力情况差异不大，在下面的抽样方法中，最合理的抽样方法是（\quad）

\choices
    {简单随机抽样}
    {按性别分层抽样}
    {按学段分层抽样}
    {系统抽样}
\end{problem}

\begin{answer}
C
\end{answer}

\begin{solution}
由于小学、初中、高中三个学段学生的视力情况差异较大，而男女生之间差异不大，应当按照差异明显的学段划分层次，再从各层中随机抽取学生. 故最合理的方法是按学段分层抽样，选 C.
\end{solution}

\begin{problem}
已知双曲线 \(C: \frac{x^2}{a^2} - \frac{y^2}{b^2} = 1\)（\(a > 0\)，\(b > 0\)）的离心率为 \(\frac{\sqrt{5}}{2}\)，则 \(C\) 的渐近线方程为（\quad）

\choices
    {\(y = \pm \frac{1}{4} x\)}
    {\(y = \pm \frac{1}{3}x\)}
    {\( y = \pm \frac{1}{2} x \)}
    {\(y = \pm x\)}
\end{problem}

\begin{answer}
C
\end{answer}

\begin{solution}
双曲线的离心率满足 \(e=\dfrac ca\)，且 \(c^2=a^2+b^2\). 由
\(\dfrac{\sqrt{5}}2=\dfrac ca\) 得 \(c^2=\dfrac54a^2\)，从而 \(b^2=c^2-a^2=\dfrac14a^2\)，即 \(\dfrac ba=\dfrac12\). 双曲线的渐近线为 \(y=\pm\dfrac ba x\)，所以方程为 \(y=\pm\dfrac12x\)，选 C.
\end{solution}

\begin{problem}
执行下面的程序框图，若输入的 \(t \in \left[-1, 3\right]\)，则输出的 \(s\) 属于（\quad）

\bitmapfigure[max width=0.42\linewidth]{img/2013/new_curriculum_paper_1_science/q5_fig1.png}

\choices
    {\(\left[-3, 4\right]\)}
    {\(\left[-5, 2\right]\)}
    {\(\left[-4, 3\right]\)}
    {\(\left[-2, 5\right]\)}
\end{problem}

\begin{answer}
A
\end{answer}

\begin{solution}
当 \(t<1\) 时，程序输出 \(s=3t\)，由 \(t\in[-1,1)\) 得 \(s\in[-3,3)\). 当 \(t\geqslant1\) 时，\(s=4t-t^2=4-(t-2)^2\)，且 \(t\in[1,3]\)，故 \(s\in[3,4]\). 合并两种情况，输出范围为 \([-3,4]\)，选 A.
\end{solution}

\begin{problem}
如图，有一个水平放置的透明无盖的正方体容器，容器高 \(8\mathrm{~cm}\)，将一个球放在容器口，再向容器内注水，当球面恰好接触水面时测得水深为 \(6\mathrm{~cm}\)，如果不计容器的厚度，则球的体积为（\quad）

\begin{minipage}{\linewidth}
\bitmapfigure[max width=0.42\linewidth]{img/2013/new_curriculum_paper_1_science/q6_fig1.png}
\end{minipage}

\choices
    {\(\frac{500\pi}{3}\mathrm{~cm}^{3}\)}
    {\(\frac{866\pi}{3}\mathrm{~cm}^{3}\)}
    {\(\frac{1372\pi}{3}\mathrm{~cm}^{3}\)}
    {\(\frac{2048\pi}{3}\mathrm{~cm}^{3}\)}
\end{problem}

\begin{answer}
A
\end{answer}

\begin{solution}
设球的半径为 \(R\). 容器口是边长为 \(8\mathrm{~cm}\) 的正方形，球放在容器口时，球面在容器口所在水平面内的截面圆内切于该正方形，因此截面圆半径为 \(4\mathrm{~cm}\). 水面在距容器底 \(6\mathrm{~cm}\) 处，且球面恰好接触水面，所以球心到水面的距离为 \(R\)，球心到容器口所在平面的距离为 \(R-2\). 在过球心且垂直于容器口的截面内，由勾股定理得
\(R^2=(R-2)^2+4^2\)，解得 \(R=5\). 因而球的体积为 \(\dfrac43\pi R^3=\dfrac43\pi\cdot5^3=\dfrac{500\pi}{3}\)，选 A.
\end{solution}

\begin{problem}
设等差数列 \(\{a_{n}\}\) 的前 \(n\) 项和为 \(S_{n}\)，\(S_{m-1} = -2\)，\(S_{m} = 0\)，\(S_{m+1} = 3\)，则 \(m =\)（\quad）

\choices
    {3}
    {4}
    {5}
    {6}
\end{problem}

\begin{answer}
C
\end{answer}

\begin{solution}
由 \(a_m=S_m-S_{m-1}=2\)，\(a_{m+1}=S_{m+1}-S_m=3\)，得等差数列的公差 \(d=1\). 又 \(S_m=0\)，所以
\(0=S_m=\dfrac m2(a_1+a_m)=\dfrac m2(a_1+2)\)，从而 \(a_1=-2\). 因此 \(a_m=a_1+(m-1)d=-2+m-1=m-3\). 由 \(a_m=2\) 得 \(m=5\)，选 C.
\end{solution}

\begin{problem}
某几何体的三视图如图所示，则该几何的体积为（\quad）

\bitmapfigure[max width=0.62\linewidth]{img/2013/new_curriculum_paper_1_science/q8_fig1.png}

\choices
    {\(16 + 8\pi\)}
    {\(8 + 8\pi\)}
    {\(16 + 16\pi\)}
    {\(8 + 16\pi\)}
\end{problem}

\begin{answer}
A
\end{answer}

\begin{solution}
由三视图可知，该几何体由一个长、宽、高分别为 \(4\)，\(2\)，\(2\) 的长方体和一个底面半径为 \(2\)、高为 \(4\) 的半圆柱组成. 因此其体积为
\(V=4\cdot2\cdot2+\dfrac12\pi\cdot2^2\cdot4=16+8\pi\)，选 A.
\end{solution}

\begin{problem}
设 \(m\) 为正整数，\(\left(x + y\right)^{2m}\) 展开式的二项式系数的最大值为 \(a\)，\(\left(x + y\right)^{2m + 1}\) 展开式的二项式系数的最大值为 \(b\)，若 \(13a = 7b\)，则 \(m =\)（\quad）

\choices
    {5}
    {6}
    {7}
    {8}
\end{problem}

\begin{answer}
B
\end{answer}

\begin{solution}
在 \((x+y)^{2m}\) 的展开式中，最大的二项式系数为 \(a=\mathrm C_{2m}^{m}\). 在 \((x+y)^{2m+1}\) 的展开式中，最大的二项式系数为 \(b=\mathrm C_{2m+1}^{m}=\mathrm C_{2m+1}^{m+1}\). 因而
\(b=\dfrac{2m+1}{m+1}\mathrm C_{2m}^{m}=\dfrac{2m+1}{m+1}a\). 由 \(13a=7b\) 得 \(13=\dfrac{7(2m+1)}{m+1}\)，即 \(13m+13=14m+7\)，所以 \(m=6\)，选 B.
\end{solution}

\begin{problem}
已知椭圆 \(E: \frac{x^2}{a^2} + \frac{y^2}{b^2} = 1\)（\(a > b > 0\)）的右焦点为 \(F\left(3,0\right)\)，过点 \(F\) 的直线交 \(E\) 于 \(A\)，\(B\) 两点. 若 \(AB\) 的中点坐标为 \(\left(1, -1\right)\)，则 \(E\) 的方程为（\quad）

\choices
    {\(\frac{x^2}{45} + \frac{y^2}{36} = 1\)}
    {\(\frac{x^2}{36} + \frac{y^2}{27} = 1\)}
    {\(\frac{x^2}{27} + \frac{y^2}{18} = 1\)}
    {\(\frac{x^2}{18} + \frac{y^2}{9} = 1\)}
\end{problem}

\begin{answer}
D
\end{answer}

\begin{solution}
焦点 \(F=(3,0)\)，弦 \(AB\) 的中点为 \(M=(1,-1)\)，所以直线 \(AB\) 的方向向量可取 \(\bs d=(2,1)\). 设弦的两个端点为 \(M+\lambda\bs d\) 和 \(M-\lambda\bs d\). 将这两个点分别代入椭圆方程并相减，得到
\(\dfrac{8\lambda}{a^2}-\dfrac{4\lambda}{b^2}=0\). 因为 \(\lambda\neq0\)，所以 \(\dfrac{2}{a^2}=\dfrac{1}{b^2}\)，即 \(a^2=2b^2\). 又右焦点横坐标为 \(3\)，故 \(c=3\)，满足 \(a^2-b^2=c^2=9\). 解得 \(b^2=9\)，\(a^2=18\)，所以椭圆方程为 \(\dfrac{x^2}{18}+\dfrac{y^2}{9}=1\)，选 D.
\end{solution}

\begin{problem}
已知函数 \( f(x) = \begin{cases}
-x^2 + 2x, & x \leqslant 0,\\
\ln\left(x + 1\right), & x > 0,
\end{cases} \) 若 \( |f(x)| \geqslant ax \)，则 \( a \) 的取值范围是（\quad）

\choices
    {\(\left(-\infty, 0\right]\)}
    {\(\left(-\infty, 1\right]\)}
    {\(\left[-2, 1\right]\)}
    {\(\left[-2, 0\right]\)}
\end{problem}

\begin{answer}
D
\end{answer}

\begin{solution}
当 \(x\leqslant0\) 时，\(f(x)\leqslant0\)，所以 \(\lvert f(x)\rvert=x^2-2x\). 不等式化为
\(x^2-2x\geqslant ax\)，即 \(x(x-a-2)\geqslant0\). 对所有 \(x\leqslant0\) 成立的充要条件是 \(a\geqslant-2\).

当 \(x>0\) 时，\(f(x)=\ln(1+x)>0\)，不等式化为 \(\dfrac{\ln(1+x)}x\geqslant a\). 由于 \(\dfrac{\ln(1+x)}x\to0\)（\(x\to+\infty\)），若 \(a>0\)，则对充分大的 \(x\) 有 \(\dfrac{\ln(1+x)}x<a\)，不等式不成立；若 \(a\leqslant0\)，则 \(ax\leqslant0<\ln(1+x)\)，不等式对所有 \(x>0\) 成立. 因而此时 \(a\leqslant0\)，最终 \(a\in[-2,0]\)，选 D.
\end{solution}

\begin{problem}
设 \(\triangle A_{n}B_{n}C_{n}\) 的三边长分别为 \(a_{n}\)，\(b_{n}\)，\(c_{n}\)，\(\triangle A_{n}B_{n}C_{n}\) 的面积为 \(S_{n}\)，\(n = 1\)，\(2\)，\(3\)，\(\cdots\)，若 \(b_{1} > c_{1}\)，\(b_{1} + c_{1} = 2a_{1}\)，\(a_{n+1} = a_{n}\)，\(b_{n+1} = \frac{c_{n} + a_{n}}{2}\)，\(c_{n+1} = \frac{b_{n} + a_{n}}{2}\)，则（\quad）

\choices
    {\(\{S_{n}\}\) 为递减数列}
    {\(\{S_{n}\}\) 为递增数列}
    {\(\{S_{2n - 1}\}\) 为递增数列，\(\{S_{2n}\}\) 为递减数列}
    {\(\{S_{2n - 1}\}\) 为递减数列，\(\{S_{2n}\}\) 为递增数列}
\end{problem}

\begin{answer}
B
\end{answer}

\begin{solution}
记 \(u_n=b_n-a_n\)，\(v_n=c_n-a_n\). 由 \(a_{n+1}=a_n\) 及 \(b_1+c_1=2a_1\)，可知 \(a_n=a_1\)，且 \(u_1+v_1=0\). 递推关系给出
\(u_{n+1}=\dfrac{v_n}{2}\)，\(v_{n+1}=\dfrac{u_n}{2}\)，并且 \(u_{n+1}+v_{n+1}=\frac{u_n+v_n}{2}\).
由 \(b_1>c_1\) 及 \(b_1+c_1=2a_1\) 得 \(u_1=b_1-a_1>0\). 由 \(u_1+v_1=0\) 归纳得 \(u_n+v_n=0\) 对所有 \(n\) 成立，所以 \(v_n=-u_n\)，且 \(u_{n+1}=-\dfrac12u_n\). 因此 \(u_n\neq0\) 对所有 \(n\) 成立，从而
\(u_{n+1}^2=\frac{v_n^2}{4}=\frac{u_n^2}{4}\).
又 \(b_n+c_n=2a_n=2a_1\)，半周长 \(p_n=\dfrac{a_n+b_n+c_n}{2}=\dfrac{3a_1}{2}\) 及 \(p_n-a_n=\dfrac{a_1}{2}\) 均为常数. 由于每个 \(\triangle A_nB_nC_n\) 都存在，\(S_n>0\).

由海伦公式，\(S_n^2=p_n(p_n-a_n)(p_n-b_n)(p_n-c_n)\). 因为 \(b_n=a_1+u_n\)，\(c_n=a_1-u_n\)，所以
\((p_n-b_n)(p_n-c_n)=\left(\dfrac{a_1}{2}-u_n\right)\left(\dfrac{a_1}{2}+u_n\right)=\dfrac{a_1^2}{4}-u_n^2\). 由于 \(u_{n+1}^2=\dfrac14u_n^2<u_n^2\)，有 \(S_{n+1}^2>S_n^2\)，从而 \(S_{n+1}>S_n\). 故 \(\{S_n\}\) 为递增数列，选 B.
\end{solution}

\section{填空题}

\begin{problem}
已知两个单位向量 \(\bs{a}\)，\(\bs{b}\) 的夹角为 \(60^{\circ}\)，\(\bs{c} = t\bs{a} + (1 - t)\bs{b}\)，若 \(\bs{b} \cdot \bs{c} = 0\)，则 \(t =\fillinblank{}\).
\end{problem}

\begin{answer}
\(2\)
\end{answer}

\begin{solution}
由 \(\lvert\bs a\rvert=\lvert\bs b\rvert=1\) 及夹角为 \(60^\circ\)，得 \(\bs b\cdot\bs a=\cos60^\circ=\dfrac12\). 因而
\(0=\bs b\cdot\bs c=t(\bs b\cdot\bs a)+(1-t)(\bs b\cdot\bs b)=\dfrac t2+1-t=1-\dfrac t2\)，解得 \(t=2\).
\end{solution}

\begin{problem}
若数列 \(\{a_{n}\}\) 的前 \(n\) 项和为 \(S_{n}=\frac{2}{3}a_{n}+\frac{1}{3}\)，则数列 \(\{a_{n}\}\) 的通项公式是 \(a_{n}=\)\fillinblank{}.
\end{problem}

\begin{answer}
\((-2)^{n-1}\)
\end{answer}

\begin{solution}
当 \(n=1\) 时，\(S_1=a_1\)，所以 \(a_1=\dfrac23a_1+\dfrac13\)，解得 \(a_1=1\). 当 \(n\geqslant2\) 时，由 \(S_n-S_{n-1}=a_n\) 得
\(a_n=\dfrac23(a_n-a_{n-1})\)，即 \(a_n=-2a_{n-1}\). 因此 \(a_n=(-2)^{n-1}\).
\end{solution}

\begin{problem}
设当 \(x = \theta\) 时，函数 \(f(x) = \sin x - 2\cos x\) 取得最大值，则 \(\cos \theta =\fillinblank{}\).
\end{problem}

\begin{answer}
\(-\frac{2\sqrt{5}}{5}\)
\end{answer}

\begin{solution}
由柯西不等式，\(f(x)=\sin x-2\cos x\leqslant\sqrt{\sin^2x+\cos^2x}\sqrt{1^2+(-2)^2}=\sqrt5\). 取最大值时等号成立，所以 \((\sin\theta,\cos\theta)\) 与 \((1,-2)\) 同向，且两者长度分别为 \(1\) 和 \(\sqrt5\). 因而 \(\cos\theta=-\dfrac2{\sqrt5}=-\dfrac{2\sqrt5}{5}\).
\end{solution}

\begin{problem}
若函数 \( f(x) = \left(1 - x^2\right)\left(x^2 + ax + b\right) \) 的图象关于直线 \( x = -2 \) 对称，则 \( f(x) \) 的最大值是\fillinblank{}.
\end{problem}

\begin{answer}
\(16\)
\end{answer}

\begin{solution}
图象关于直线 \(x=-2\) 对称，反射变换为 \(x\mapsto-4-x\). 因为 \(f(1)=f(-1)=0\)，对称性要求反射后的点 \(-5\) 和 \(-3\) 也是零点. 于是 \(x^2+ax+b=(x+3)(x+5)=x^2+8x+15\)，从而
\(f(x)=(1-x^2)(x+3)(x+5)\). 令 \(u=x+2\)，则 \(f(x)=(-u^2+4u-3)(u^2+4u+3)=-u^4+10u^2-9=16-(u^2-5)^2\leqslant16\). 当 \(u^2=5\) 时等号成立，所以 \(f(x)\) 的最大值为 \(16\).
\end{solution}

\section{解答题}

\begin{problem}
如图，在 \(\triangle ABC\) 中，\(\angle ABC = 90^{\circ}\)，\(AB = \sqrt{3}\)，\(BC = 1\)，\(P\) 为 \(\triangle ABC\) 内一点，\(\angle BPC = 90^{\circ}\).

\begin{enumerate}
\item 若 \(PB = \frac{1}{2}\)，求 \(PA\)；
\item 若 \(\angle APB = 150^{\circ}\)，求 \(\tan \angle PBA\).
\end{enumerate}

\bitmapfigure[max width=0.34\linewidth]{img/2013/new_curriculum_paper_1_science/q17_fig1.png}
\end{problem}

\begin{solution}
\begin{enumerate}
\item 取 \(B=(0,0)\)，\(A=(\sqrt3,0)\)，\(C=(0,1)\)，设 \(P=(u,v)\). 因为 \(P\) 在三角形内部，所以 \(u>0\)，\(v>0\). 由 \(\angle BPC=90^\circ\)，有
\((-u,-v)\cdot(-u,1-v)=0\)，即 \(u^2+v^2-v=0\). 若 \(PB=\dfrac12\)，则 \(u^2+v^2=\dfrac14\)，与上式联立得 \(v=\dfrac14\)，\(u^2=\dfrac{3}{16}\). 由 \(u>0\) 得 \(u=\dfrac{\sqrt3}{4}\)，所以
\(PA^2=\left(\sqrt3-\dfrac{\sqrt3}{4}\right)^2+\left(\dfrac14\right)^2=\dfrac{27}{16}+\dfrac1{16}=\dfrac74\)，故 \(PA=\dfrac{\sqrt7}{2}\).

\item 设 \(\alpha=\angle PBA\)，\(r=PB\). 由上面的坐标关系，若 \(P=(r\cos\alpha,r\sin\alpha)\)，则 \(u^2+v^2=v\) 化为 \(r^2=r\sin\alpha\). 因 \(r>0\)，所以 \(r=\sin\alpha\). 在三角形 \(ABP\) 中，\(\angle APB=150^\circ\)，且 \(P\) 在三角形内部，所以 \(0<\alpha<30^\circ\)，从而 \(\angle PAB=30^\circ-\alpha\). 由正弦定理，
\(\dfrac{PB}{\sin(30^\circ-\alpha)}=\dfrac{AB}{\sin150^\circ}=2\sqrt3\). 代入 \(PB=\sin\alpha\)，得
\(\sin\alpha=2\sqrt3\sin(30^\circ-\alpha)=\sqrt3\cos\alpha-3\sin\alpha\). 因此 \(4\sin\alpha=\sqrt3\cos\alpha\)，从而 \(\tan\angle PBA=\tan\alpha=\dfrac{\sqrt3}{4}\).
\end{enumerate}
\end{solution}

\begin{problem}
如图，三棱柱 \(ABC - A_{1}B_{1}C_{1}\) 中，\(CA = CB\)，\(AB = AA_{1}\)，\(\angle BAA_{1} = 60^{\circ}\).

\begin{enumerate}
\item 证明： \(AB \perp A_{1}C\)；
\item 若平面 \( ABC \perp \) 平面 \(AA_{1}B_{1}B\)，\( AB = CB \)，求直线 \(A_{1}C\) 与平面 \(BB_{1}C_{1}C\) 所成角的正弦值.
\end{enumerate}

\begin{center}
\bitmapfigure[float=inline,max width=\linewidth]{img/2013/new_curriculum_paper_1_science/q18_fig1.png}
\end{center}
\end{problem}

\begin{solution}
\begin{enumerate}
\item 设 \(\bs u=\overrightarrow{AB}\)，\(\bs v=\overrightarrow{AC}\)，\(\bs w=\overrightarrow{AA_1}\). 由 \(CA=CB\)，得
\(\lvert\bs v\rvert^2=\lvert\bs v-\bs u\rvert^2\)，从而 \(\bs u\cdot\bs v=\dfrac12\lvert\bs u\rvert^2\). 又 \(AB=AA_1\)，且 \(\angle BAA_1=60^\circ\)，所以
\(\bs u\cdot\bs w=\lvert\bs u\rvert\lvert\bs w\rvert\cos60^\circ=\dfrac12\lvert\bs u\rvert^2\). 因此
\(\overrightarrow{AB}\cdot\overrightarrow{A_1C}=\bs u\cdot(\bs v-\bs w)=0\)，故 \(AB\perp A_1C\).

\item 设 \(AB=CB=l\). 由 \(CA=CB\) 可知 \(CA=AB=BC=l\)，所以底面三角形 \(ABC\) 为正三角形. 取 \(A=(0,0,0)\)，\(B=(l,0,0)\)，\(C=\left(\dfrac l2,\dfrac{\sqrt3l}{2},0\right)\). 因为平面 \(ABC\perp\) 平面 \(AA_1B_1B\)，且 \(AA_1\) 与 \(AB\) 成 \(60^\circ\)，可取
\(A_1=\left(\dfrac l2,0,\dfrac{\sqrt3l}{2}\right)\). 于是 \(A_1C\) 的方向向量可取
\(\bs d=\overrightarrow{A_1C}=\left(0,\dfrac{\sqrt3l}{2},-\dfrac{\sqrt3l}{2}\right)\).

平面 \(BB_1C_1C\) 内有两个不共线方向向量
\(\overrightarrow{BC}=\left(-\dfrac l2,\dfrac{\sqrt3l}{2},0\right)\)，\(\overrightarrow{BB_1}=\overrightarrow{AA_1}=\left(\dfrac l2,0,\dfrac{\sqrt3l}{2}\right)\). 取其法向量
\(\bs n=\overrightarrow{BC}\times\overrightarrow{BB_1}=\left(\dfrac{3l^2}{4},\dfrac{\sqrt3l^2}{4},-\dfrac{\sqrt3l^2}{4}\right)\). 因而
\(\lvert\bs d\rvert=l\sqrt{\dfrac32}\)，\(\lvert\bs n\rvert=\dfrac{l^2\sqrt{15}}4\)，且 \(\lvert\bs d\cdot\bs n\rvert=\dfrac{3l^3}{4}\). 设直线 \(A_1C\) 与平面的夹角为 \(\varphi\)，则
\(\sin\varphi=\dfrac{\lvert\bs d\cdot\bs n\rvert}{\lvert\bs d\rvert\lvert\bs n\rvert}=\dfrac{3l^3/4}{l\sqrt{3/2}\cdot l^2\sqrt{15}/4}=\dfrac{2}{\sqrt{10}}=\dfrac{\sqrt{10}}5\).
\end{enumerate}
\end{solution}

\begin{problem}
一批产品需要进行质量检验，检验方案是：先从这批产品中任取 \(4\) 件作检验，这 \(4\) 件产品中优质品的件数记为 \(n\). 如果 \(n = 3\)，再从这批产品中任取 \(4\) 件作检验，若都为优质品，则这批产品通过检验；如果 \(n = 4\)，再从这批产品中任取 \(1\) 件作检验，若为优质品，则这批产品通过检验；其他情况下，这批产品都不能通过检验. 假设这批产品的优质品率为 \(50\%\). 即取出的产品是优质品的概率都为 \(\frac{1}{2}\)，且各件产品是否为优质品相互独立.
\begin{enumerate}
\item 求这批产品通过检验的概率；
\item 已知每件产品检验费用为 \(100\text{ 元}\)，凡抽取的每件产品都需要检验，对这批产品作质量检验所需的费用记为 \(X\)（单位：元），求 \(X\) 的分布列及数学期望.
\end{enumerate}
\end{problem}

\begin{solution}
\begin{enumerate}
\item 记第一次抽取的优质品件数为 \(n\). 由题意，\(P(n=3)=\mathrm C_4^3\left(\dfrac12\right)^4=\dfrac14\)，\(P(n=4)=\left(\dfrac12\right)^4=\dfrac1{16}\). 当 \(n=3\) 时，第二次抽取的 \(4\) 件全为优质品的概率为 \(\left(\dfrac12\right)^4=\dfrac1{16}\)；当 \(n=4\) 时，再抽取 \(1\) 件为优质品的概率为 \(\dfrac12\). 因此通过检验的概率为
\(P=\dfrac14\cdot\dfrac1{16}+\dfrac1{16}\cdot\dfrac12=\dfrac1{64}+\dfrac2{64}=\dfrac3{64}\).

\item 若 \(n=0\)，\(1\)，\(2\)，检验在第一次抽取后结束，共检验 \(4\) 件，所以 \(X=400\)，其概率为
\(P(X=400)=P(n\leqslant2)=1-\dfrac14-\dfrac1{16}=\dfrac{11}{16}\). 若 \(n=4\)，还需检验 \(1\) 件，所以 \(X=500\)，且 \(P(X=500)=\dfrac1{16}\). 若 \(n=3\)，还需检验 \(4\) 件，所以 \(X=800\)，且 \(P(X=800)=\dfrac14\).

故 \(X\) 的分布列为
\begin{center}
\begin{tabular}{c|ccc}
\(X\) & \(400\) & \(500\) & \(800\) \\
\hline
\(P\) & \(\dfrac{11}{16}\) & \(\dfrac1{16}\) & \(\dfrac14\)
\end{tabular}
\end{center}

数学期望为 \(E(X)=400\cdot\dfrac{11}{16}+500\cdot\dfrac1{16}+800\cdot\dfrac14=506.25\).
\end{enumerate}
\end{solution}

\begin{problem}
已知圆 \(M: \left(x + 1\right)^2 + y^2 = 1\)，圆 \(N: \left(x - 1\right)^2 + y^2 = 9\)，动圆 \(P\) 与圆 \(M\) 外切并与圆 \(N\) 内切，圆心 \(P\) 的轨迹为曲线 \(C\).
\begin{enumerate}
\item 求 \(C\) 的方程；
\item \(l\) 是与圆 \(P\)，圆 \(M\) 都相切的一条直线，\(l\) 与曲线 \(C\) 交于 \(A\)，\(B\) 两点，当圆 \(P\) 的半径最长时，求 \(|AB|\).
\end{enumerate}
\end{problem}

\begin{solution}
\begin{enumerate}
\item 设圆 \(P\) 的半径为 \(r\)，圆心坐标为 \(P(x,y)\). 圆 \(M\) 和圆 \(N\) 的圆心分别为 \((-1,0)\) 和 \((1,0)\)，半径分别为 \(1\) 和 \(3\). 由外切、内切条件，得
\(\sqrt{(x+1)^2+y^2}=r+1\)，\(\sqrt{(x-1)^2+y^2}=3-r\). 两式相加可得
\(\sqrt{(x+1)^2+y^2}+\sqrt{(x-1)^2+y^2}=4\)，所以 \(C\) 是以 \((-1,0)\)，\((1,0)\) 为焦点、长轴长为 \(4\) 的椭圆，即
\(C:\dfrac{x^2}{4}+\dfrac{y^2}{3}=1\). 由 \(r=\sqrt{(x+1)^2+y^2}-1>0\)，椭圆左端点 \((-2,0)\) 对应 \(r=0\) 不表示圆，故 \(x\neq-2\).

\item 由 \(r=\sqrt{(x+1)^2+y^2}-1\)，并利用 \(y^2=3-\dfrac34x^2\)，得
\(r+1=\sqrt{(x+1)^2+y^2}=\sqrt{\dfrac{(x+4)^2}{4}}\). 椭圆上 \(-2\leqslant x\leqslant2\)，所以 \(r\) 在 \(x=2\) 处取得最大值 \(2\)，此时圆 \(P\) 的圆心为 \((2,0)\)，半径为 \(2\).

此时圆 \(M\) 与圆 \(P\) 的内公切线为 \(x=0\)，它与椭圆的交点为 \((0,\sqrt3)\)，\((0,-\sqrt3)\)，故 \(|AB|=2\sqrt3\).

其余两条公切线为外公切线. 设其中一条为 \(y=mx+b\)，记 \(\delta=\sqrt{1+m^2}\). 两圆心在直线同侧时，点 \((-1,0)\) 和 \((2,0)\) 到直线的有向距离分别为半径 \(1\) 和 \(2\)，故可取
\(b-m=\delta\)，\(b+2m=2\delta\). 两式相减得 \(3m=\delta\)，从而 \(m^2=\dfrac18\)，取其中一条直线可写为 \(y=\dfrac{x}{2\sqrt2}+\sqrt2\). 将其代入椭圆方程，得到
\(\dfrac{x^2}{4}+\dfrac{\left(\dfrac{x}{2\sqrt2}+\sqrt2\right)^2}{3}=1\)，即 \(7x^2+8x-8=0\). 两交点的横坐标之差为 \(\dfrac{12\sqrt2}{7}\)，而直线斜率的长度因子为 \(\sqrt{1+m^2}=\dfrac{3}{2\sqrt2}\)，所以
\(|AB|=\dfrac{3}{2\sqrt2}\cdot\dfrac{12\sqrt2}{7}=\dfrac{18}{7}\). 因而 \(|AB|=2\sqrt3\) 或 \(\dfrac{18}{7}\).
\end{enumerate}
\end{solution}

\begin{problem}
设函数 \(f(x) = x^2 + ax + b\)，\(g(x) = \e^x\left(cx + d\right)\)，若曲线 \(y = f(x)\) 和曲线 \(y = g(x)\) 都过点 \(P\left(0,2\right)\)，且在点 \(P\) 处有相同的切线 \(y = 4x + 2\).
\begin{enumerate}
\item 求 \(a\)，\(b\)，\(c\)，\(d\) 的值；
\item 若 \(x \geqslant -2\) 时，\(f(x) \leqslant kg(x)\)，求 \(k\) 的取值范围.
\end{enumerate}
\end{problem}

\begin{solution}
\begin{enumerate}
\item 由 \(f(0)=2\)，得 \(b=2\). 又 \(f'(x)=2x+a\)，而点 \(P\) 处切线斜率为 \(4\)，所以 \(a=f'(0)=4\). 由 \(g(0)=2\) 得 \(d=2\). 又
\(g'(x)=\e^x(cx+d+c)\)，故 \(g'(0)=c+d=4\)，从而 \(c=2\). 因此 \(a=4\)，\(b=2\)，\(c=2\)，\(d=2\).

\item 代入所得系数，\(f(x)=x^2+4x+2\)，\(g(x)=2\e^x(x+1)\). 在 \(x=-1\) 时，\(f(-1)=-1\)，\(g(-1)=0\)，所以不等式成立. 对 \(x\neq-1\)，令
\(R(x)=\dfrac{f(x)}{g(x)}=\dfrac{\e^{-x}}2\left(x+3-\dfrac1{x+1}\right)\).

当 \(-2\leqslant x<-1\) 时，\(g(x)<0\)，不等式 \(f(x)\leqslant kg(x)\) 等价于 \(k\leqslant R(x)\). 求导得
\(R'(x)=\dfrac{\e^{-x}}2\dfrac{(x+2)^2(-x)}{(x+1)^2}\geqslant0\)，所以 \(R(x)\) 在 \([-2,-1)\) 上递增，且 \(R(-2)=\e^2\)，从而此区间要求 \(k\leqslant\e^2\).

当 \(x>-1\) 时，\(g(x)>0\)，不等式等价于 \(k\geqslant R(x)\). 同样由导数符号知 \(R(x)\) 在 \((-1,0)\) 上递增，在 \((0,+\infty)\) 上递减，所以最大值为 \(R(0)=1\)，从而要求 \(k\geqslant1\). 综合得 \(k\in[1,\e^2]\).
\end{enumerate}
\end{solution}

\begin{problem}
如图，直线 \(AB\) 为圆的切线，切点为 \(B\)，点 \(C\) 在圆上，\(\angle ABC\) 的角平分线 \(BE\) 交圆于点 \(E\)，\(DB\) 垂直 \(BE\) 交圆于 \(D\).

\begin{enumerate}
\item 证明： \(DB = DC\)；
\item 设圆的半径为 \(1\)，\(BC = \sqrt{3}\)，延长 \(CE\) 交 \(AB\) 于点 \(F\)，求 \(\triangle BCF\) 外接圆的半径.
\end{enumerate}

\bitmapfigure[max width=0.34\linewidth]{img/2013/new_curriculum_paper_1_science/q22_fig1.png}
\end{problem}

\begin{solution}
\begin{enumerate}
\item 由弦切角定理，\(\angle ABE=\angle BCE\)；又因为 \(BE\) 平分 \(\angle ABC\)，所以 \(\angle ABE=\angle EBC\). 因而 \(\angle BCE=\angle EBC\)，从而 \(BE=CE\). 由于 \(DB\perp BE\)，有 \(\angle DBE=90^\circ\)，所以 \(DE\) 是圆的直径，进而 \(\angle DCE=90^\circ\). 在直角三角形 \(DBE\) 和 \(DCE\) 中，它们有公共斜边 \(DE\)，且 \(BE=CE\)，故两三角形全等，从而 \(DB=DC\).

\item 建立直角坐标系，使 \(B=(0,0)\)，直线 \(AB\) 为 \(x\) 轴正半轴，圆心为 \(O=(0,1)\). 因圆的半径为 \(1\)，圆的方程为 \(x^2+(y-1)^2=1\). 由 \(BC=\sqrt3\) 及 \(C\) 在圆上，得 \(x_C^2+y_C^2=3\) 且 \(x_C^2+(y_C-1)^2=1\)，两式相减得 \(y_C=\dfrac32\)；按图取 \(C=\left(\dfrac{\sqrt3}{2},\dfrac32\right)\)，所以 \(\angle ABC=60^\circ\).

射线 \(BE\) 是 \(\angle ABC\) 的角平分线，所以可取其方向向量为 \(\left(\dfrac{\sqrt3}{2},\dfrac12\right)\). 令 \(E=t\left(\dfrac{\sqrt3}{2},\dfrac12\right)\)，代入圆方程 \(x^2+y^2-2y=0\)，得 \(t^2-t=0\). 除去 \(t=0\) 对应的 \(B\)，有 \(t=1\)，故 \(E=\left(\dfrac{\sqrt3}{2},\dfrac12\right)\).

直线 \(DB\perp BE\)，可取其方向向量为 \(\left(-\dfrac12,\dfrac{\sqrt3}{2}\right)\). 令 \(D=t\left(-\dfrac12,\dfrac{\sqrt3}{2}\right)\)，代入圆方程得 \(t^2-\sqrt3t=0\). 除去 \(t=0\) 后 \(t=\sqrt3\)，所以 \(D=\left(-\dfrac{\sqrt3}{2},\dfrac32\right)\). 于是 \(DB=\sqrt3\)，且
\(DC=\sqrt{\left(\sqrt3\right)^2+0^2}=\sqrt3\)，故 \(DB=DC\).

由 \(C\) 与 \(E\) 的横坐标相同，直线 \(CE\) 为 \(x=\dfrac{\sqrt3}{2}\)，它与 \(AB\) 的交点为 \(F=\left(\dfrac{\sqrt3}{2},0\right)\). 因此 \(BF=\dfrac{\sqrt3}{2}\)，\(CF=\dfrac32\)，且 \(\angle BFC=90^\circ\). 三角形 \(BCF\) 的斜边为 \(BC=\sqrt3\)，所以其外接圆半径为斜边的一半，即 \(\dfrac{\sqrt3}{2}\).
\end{enumerate}
\end{solution}

\begin{problem}
已知曲线 \(C_1\) 的参数方程为 \(\begin{cases}
x = 4 + 5\cos t,\\
y = 5 + 5\sin t,
\end{cases}\) （\(t\) 为参数），以坐标原点为极点，\(x\) 轴的正半轴为极轴建立极坐标系. 曲线 \(C_2\) 的极坐标方程为 \(\rho = 2\sin \theta\).
\begin{enumerate}
\item 把 \(C_1\) 的参数方程化为极坐标方程；
\item 求 \(C_1\) 与 \(C_2\) 交点的极坐标 \(\left(\rho \geqslant 0, 0 \leqslant \theta < 2\pi\right)\).
\end{enumerate}
\end{problem}

\begin{solution}
\begin{enumerate}
\item 由参数方程知 \(C_1\) 是圆 \((x-4)^2+(y-5)^2=25\). 代入 \(x=\rho\cos\theta\)，\(y=\rho\sin\theta\)，得
\(\rho^2-8\rho\cos\theta-10\rho\sin\theta+16=0\).

\item 将 \(\rho=2\sin\theta\) 代入上式，得
\(4\sin^2\theta-16\sin\theta\cos\theta-20\sin^2\theta+16=0\)，即
\(1-\sin^2\theta-\sin\theta\cos\theta=0\). 化简为
\(\cos\theta(\cos\theta-\sin\theta)=0\).

因为 \(\rho\geqslant0\) 且 \(\rho=2\sin\theta\)，所以 \(\sin\theta\geqslant0\)，在 \(0\leqslant\theta<2\pi\) 内有 \(0\leqslant\theta\leqslant\pi\). 当 \(\cos\theta=0\) 时，\(\theta=\dfrac\pi2\)，\(\rho=2\)；当 \(\cos\theta=\sin\theta\) 时，\(\theta=\dfrac\pi4\)，\(\rho=\sqrt2\). 因此交点的极坐标为 \(\left(\sqrt2,\dfrac\pi4\right)\) 和 \(\left(2,\dfrac\pi2\right)\).
\end{enumerate}
\end{solution}

\begin{problem}
已知函数 \( f(x) = |2x - 1| + |2x + a| \)，\( g(x) = x + 3 \).
\begin{enumerate}
\item 当 \(a = -2\) 时，求不等式 \(f(x) < g(x)\) 的解集；
\item 设 \(a > -1\)，且当 \(x \in \left[-\frac{a}{2}, \frac{1}{2}\right]\) 时，\(f(x) \leqslant g(x)\)，求 \(a\) 的取值范围.
\end{enumerate}
\end{problem}

\begin{solution}
\begin{enumerate}
\item 当 \(a=-2\) 时，\(f(x)=|2x-1|+|2x-2|\). 分 \(x<\dfrac12\)、\(\dfrac12\leqslant x<1\)、\(x\geqslant1\) 三种情况讨论：
\(x<\dfrac12\) 时，\(f(x)=3-4x\)，由 \(f(x)<x+3\) 得 \(x>0\)，所以得到 \((0,\dfrac12)\)；
\(\dfrac12\leqslant x<1\) 时，\(f(x)=1\)，不等式恒成立，所以得到 \([\dfrac12,1)\)；
\(x\geqslant1\) 时，\(f(x)=4x-3\)，由 \(f(x)<x+3\) 得 \(x<2\)，所以得到 \([1,2)\). 合并得解集为 \((0,2)\).

\item 因为 \(a>-1\)，所以 \(-\dfrac a2<\dfrac12\). 当 \(x\in[-\dfrac a2,\dfrac12]\) 时，\(2x-1\leqslant0\)，\(2x+a\geqslant0\)，故
\(f(x)=1-2x+2x+a=a+1\). 要使 \(f(x)\leqslant g(x)=x+3\) 对该区间内所有 \(x\) 成立，只需在左端点检验：
\(a+1\leqslant3-\dfrac a2\)，即 \(a\leqslant\dfrac43\). 与 \(a>-1\) 合并，得 \(a\in\left(-1,\dfrac43\right]\).
\end{enumerate}
\end{solution}
